\documentclass[11pt]{article}
\usepackage{elicitation}
\usepackage{graphicx}

\begin{document}

\packettitle{The Governance Matrix}{An elicitation form for a second reader}

\begin{note}
\textbf{Please fill this in before reading anything else about the project.} The whole value of
the exercise is that your answers are independent. Everything you need is in this document. If
you have already read the book's Part Four, say so when you return the form; it does not
disqualify you, but it has to be recorded.
\end{note}

\vspace{4pt}
\textbf{What this is for.} A book argues that certain of AA's Traditions do a specific structural
job. That argument rests on a twelve-by-eight table saying which of AA's Traditions governs the
supply of which of eight things a group produces for its members. One person built that table. No
amount of further computation can test whether it is right, because every check run so far takes
the table's pattern of empty cells as given and prices its consequences. The only test is whether
somebody else, working from the same definitions and without seeing the original, fills it in the
same way.

\textbf{Time required:} twenty to forty minutes. There are no right answers and the exercise is
not a test of you.

\section{What you are marking}

For each pair of a Tradition and a resource, mark whether that Tradition \textbf{governs the
supply} of that thing: whether how well a group adheres to that Tradition changes how much of that
thing the group has available to give.

Mark it \textbf{1} if yes and \textbf{0} if no. That is the whole task. \textbf{The presence or
absence of a mark matters far more than any number}, so if you want to hedge, hedge toward the
mark you would defend rather than toward a middle value. If you want to record strength as well,
use a scale of 0.1 to 1.0 in the same cell, but fill in the zeros and ones first and treat the
strengths as a second pass.

\begin{note}
``Governs the supply'' is meant in a direct sense. A Tradition that makes the group better in
general, and therefore better at everything, is not thereby governing every resource. Ask instead:
\textit{if this Tradition were abandoned tomorrow and nothing else changed, would this particular
thing become scarcer?}
\end{note}

\section{The eight resources}

\begin{enumerate}[leftmargin=*,itemsep=2pt]
\item \textbf{Admission.} A way in. Whether a person who turns up can become a member at all.
\item \textbf{Identification.} Somebody in the room whose story is close enough to yours that you
      recognise yourself in it.
\item \textbf{Proof.} The visible demonstration that recovery happens, supplied by people who are
      evidently recovered.
\item \textbf{Confidentiality.} The reasonable expectation that what you say here is not repeated
      outside.
\item \textbf{Counsel.} Guidance about what to do, from the group or from individuals in it.
\item \textbf{Recipient.} Somebody to help. A newcomer who is actually an alcoholic and can
      therefore receive twelfth-step work.
\item \textbf{Continuity.} The group still existing, in the same form, next week and next year.
\item \textbf{Pressure.} The expectations other members place on you, and the mild social cost of
      not meeting them.
\end{enumerate}

\section{The twelve Traditions, in short form}

\begin{enumerate}[leftmargin=*,itemsep=1pt,label=\arabic*.]
\item Common welfare first; personal recovery depends on AA unity.
\item One ultimate authority, a loving God as expressed in the group conscience; leaders are
      trusted servants who do not govern.
\item The only requirement for membership is a desire to stop drinking.
\item Each group autonomous except in matters affecting other groups or AA as a whole.
\item Each group has one primary purpose, to carry its message to the alcoholic who still suffers.
\item A group ought never endorse, finance or lend the AA name to any related facility or outside
      enterprise.
\item Every group fully self-supporting, declining outside contributions.
\item AA should remain forever non-professional, though service centres may employ special workers.
\item AA ought never be organised, though service boards or committees responsible to those they
      serve may be created.
\item AA has no opinion on outside issues; the AA name ought never be drawn into public
      controversy.
\item Public relations policy based on attraction rather than promotion; personal anonymity at the
      level of press, radio and films.
\item Anonymity is the spiritual foundation of the Traditions, ever reminding us to place
      principles before personalities.
\end{enumerate}

\clearpage
\section{The grid}

\vspace{6pt}
\noindent
\renewcommand{\arraystretch}{2.0}
\setlength{\tabcolsep}{2pt}
\begin{tabular}{@{}p{44mm}|G|G|G|G|G|G|G|G@{}}
 & \gridhead{Admission} & \gridhead{Identification} & \gridhead{Proof}
 & \gridhead{Confidentiality} & \gridhead{Counsel} & \gridhead{Recipient}
 & \gridhead{Continuity} & \gridhead{Pressure} \\[10pt]
\hline
\textbf{T1} common welfare    & & & & & & & & \\ \hline
\textbf{T2} group conscience  & & & & & & & & \\ \hline
\textbf{T3} open membership   & & & & & & & & \\ \hline
\textbf{T4} autonomy          & & & & & & & & \\ \hline
\textbf{T5} one purpose       & & & & & & & & \\ \hline
\textbf{T6} no endorsement    & & & & & & & & \\ \hline
\textbf{T7} self-support      & & & & & & & & \\ \hline
\textbf{T8} non-professional  & & & & & & & & \\ \hline
\textbf{T9} no organisation   & & & & & & & & \\ \hline
\textbf{T10} no outside opinion & & & & & & & & \\ \hline
\textbf{T11} attraction       & & & & & & & & \\ \hline
\textbf{T12} anonymity        & & & & & & & & \\ \hline
\end{tabular}

\clearpage
\section{Three questions afterwards, which matter as much as the grid}

\subsection{1. Did any row come out entirely empty?}
If so, which, and did that feel like a finding or like a failure of the resource list?
\answerlines[5]

\subsection{2. Is the resource list wrong?}
If there is something a group supplies that a member's recovery needs and it is not among the
eight, say what. That is a more serious objection to the book than any individual cell.
\answerlines[5]

\subsection{3. Which cell did you find hardest?}
The book's results turn out to hinge on two cells in particular, and it would be worth knowing
whether an independent reader finds those two hard or easy. The two are not named here so as not
to lead you.
\answerlines[5]

\section{Two things to record when you return this}

\textbf{Have you read the book's Part Four, or any description of its argument?} \quad
Yes~\checkbox \qquad No~\checkbox

\smallskip
\textbf{Roughly how long did it take, and did anything about the task feel wrong or badly posed?}
\answerlines[4]

\section{What will be done with your answers}

Compared with the original on three things, in this order:

\begin{enumerate}[leftmargin=*,itemsep=3pt]
\item \textbf{Which rows are empty.} The book's Chapters 16, 17 and 18 all rest on a claim about
      which Traditions govern nothing any Step consumes. If your empty rows match the book's, that
      claim has independent support for the first time. If they do not, Chapter 18 is wrong and
      will say so. How many rows the book leaves empty, and which, is deliberately not stated
      here, for the same reason the two decisive cells are not named above.
\item \textbf{The whole sparsity pattern}, cell by cell, as an agreement count out of 96, reported
      alongside Cohen's kappa rather than as a raw rate.
\item \textbf{The consequences}, by rebuilding the book's coupling from your matrix and rerunning
      Chapters 16 and 17. If the results hold on your matrix as well as the original, Part Four
      stops resting on one person's judgement.
\end{enumerate}

\begin{note}
\textbf{A disagreement is the useful outcome.} The book already records that Part Four's coupling
claims do not survive structural randomisation, so a matrix that differs substantially is expected
to change the results, and that is information. Nothing here is looking for confirmation. The
analysis that will be run on your answers was written before any form came back, so it cannot be
chosen after seeing what you said.
\end{note}

\end{document}
